\diarytitle{0071}{2階線形偏微分方程式の型判定と領域による型の変化（$y u_{xx} + u_{yy} = 0$）}

方針: 2階線形偏微分方程式
\[
A(x,y)\,u_{xx} + 2B(x,y)\,u_{xy} + C(x,y)\,u_{yy} + (\text{低階の項}) = 0
\]
の型は、判別式 $D = B^{2} - AC$ の符号によって次のように分類される。
\[
D = B^{2}-AC
\begin{cases}
< 0 & \text{楕円型 (elliptic)}\\
= 0 & \text{放物型 (parabolic)}\\
> 0 & \text{双曲型 (hyperbolic)}
\end{cases}
\]
係数が変数の場合、$D$ の符号が点ごとに変わりうるため、領域を分けて型を判定する必要がある。

$A=y,\ B=0,\ C=1$ であるから
\[
D = B^{2}-AC = 0 - y\cdot 1 = -y
\]
よって
\[
D = -y
\begin{cases}
<0 & (y>0) \Rightarrow \text{楕円型}\\
=0 & (y=0) \Rightarrow \text{放物型（退化）}\\
>0 & (y<0) \Rightarrow \text{双曲型}
\end{cases}
\]
（これは Tricomi 方程式と呼ばれる有名な混合型方程式であり、$y=0$ が型の変わる境界である。）
\[
\boxed{\; y>0:\text{楕円型},\quad y=0:\text{放物型},\quad y<0:\text{双曲型}\;}
\]
（検算: $y=1>0$ では $D=-1<0$ で楕円型、$y=-1<0$ では $D=1>0$ で双曲型となり、符号の対応が一致する。）
